\documentclass[11pt]{article}

\usepackage[utf8]{inputenc}
\usepackage{color}
\usepackage{graphicx}
\usepackage{geometry}
\usepackage{hyperref}
\hypersetup{
    colorlinks,
    citecolor=black,
    filecolor=black,
    linkcolor=black,
    urlcolor=black
}
\usepackage{pdfpages}

\renewcommand{\contentsname}{Inhaltsverzeichnis}

\begin{document}

\begin{titlepage}
	\begin{center}
		\vspace*{1cm}

		\Huge
		\textbf{Tourney}\\
		Systemtestplan

		\vspace{0.5cm}
		\LARGE
		Softwarepraktikum\\
		\Large
		Wintersemester 2014/2015

		\vspace{1.5cm}

		\large
		\textbf{Jonas Auer (2860992)\\
				 Fabian Biester (2859084)\\
				 Jan Tagscherer (2893134)}

		\vfill

		\includegraphics[width=0.4\textwidth]{../Logo.png}

		\vspace{1.5cm}

		\Large
		Universität Stuttgart\\
		31.12.2014
	\end{center}
\end{titlepage}

\newpage

\tableofcontents
\newpage

\section{Einleitung}

\subsection{Zweck des Systemtestplans}

Mit dem Systemtest, der in diesem Dokument geplant wird, wird die Turniermanagement-Software \textbf{Tourney} getestet.

Durch das Festhalten aller Testfälle wird ein Abschätzen des Aufwands und ein strukturiertes Testen des kompletten Systems ermöglicht. Dadurch wird die Qualität der Software und ein Abdecken aller Anwendungsfälle gewährleistet.

\newpage

\section{Vorgehen}

\subsection{Testumgebung}

Um das System als ganzes zu testen und alle Testfälle durchführen zu können muss eine Testumgebung mit den folgenden Bedingungen gegeben sein.

\begin{itemize}
	\item Ein Arbeitsplatz mit einem gängigen Betriebsssytem (Windows/Mac/Linux) und mit allen Bedingungen, um Java-Software auszuführen (Aktuelles JRE). Dieser Rechner dient als Arbeitsplatz des Administrators.
	\item Ein Arbeitsplatz mit einem gängigen Betriebsssytem (Windows/Mac/Linux) und mit allen Bedingungen, um Java-Software auszuführen (Aktuelles JRE). Dieser Rechner dient als Arbeitsplatz, auf dem zusätzlich Anmeldungen durchgeführt werden können.
	\item Ein Arbeitsplatz mit einem gängigen Betriebsssytem (Windows/Mac/Linux) und mit allen Bedingungen, um Java-Software auszuführen (Aktuelles JRE). Dieser Rechner dient als Arbeitsplatz eines Turnierordners, der ein hinzugefügtes Turnier durchführt.
	\item Ein Wechseldatenträger (USB-Stick), um Daten zwischen den einzelnen Arbeitsplätzen zu übertragen
\end{itemize}

Die Tests werden manuell durchgeführt und ausgewertet. Die Ergebnisse werden mittels der Testfallverwaltung \textit{TSM} protokolliert.

\subsection{Testzeitraum}

Der Systemtest wird im Zeitraum vom 01.01.2015 bis zum 23.01.2015 durchgeführt.

\newpage

\section{Testfälle}

\includepdf[pages={-}]{testfaelle.pdf}

\newpage

\section{Versionshistorie}

\begin{itemize}
	\item Version 1.0 (30.12.2014)
	\begin{itemize}
		\item Grundstruktur des Systemtestplans erstellt und Einleitung verfasst
	\end{itemize}
	\item Version 1.1 (31.12.2014)
	\begin{itemize}
		\item Testfälle hinzugefügt
	\end{itemize}
\end{itemize}

\end{document}
